\section{Источники погрешности в математическом моделировании}
\label{sec:error-sources-in-modeling}

Результат вычислительного моделирования является приближённым по нескольким
принципиально различным причинам. Удобно разделять ошибку по этапам
перехода
\[
\text{реальный объект}
\to
\text{математическая модель}
\to
\text{численный метод}
\to
\text{вычисления в ЭВМ}.
\]

Пусть $y_{\rm real}$ --- интересующая характеристика реального объекта,
$y$ --- точное решение принятой математической модели, $y_h$ --- точный
результат дискретного метода при параметре $h$, а $\widehat y_h$ --- реально
вычисленное в конечной арифметике значение. Тогда формально
\[
\widehat y_h-y_{\rm real}
=
\underbrace{(y-y_{\rm real})}_{\text{погрешность модели и данных}}
+
\underbrace{(y_h-y)}_{\text{погрешность метода}}
+
\underbrace{(\widehat y_h-y_h)}_{\text{вычислительная погрешность}}.
\]
Это разложение подчёркивает, что уменьшение одной составляющей не устраняет
остальные.

\subsection{Неустранимая погрешность и погрешность математической модели}

Источниками первой группы являются
\begin{itemize}
    \item идеализация реального объекта и отброшенные физические эффекты;
    \item неточно известные параметры, коэффициенты, начальные и граничные
    данные;
    \item погрешности измерений;
    \item неполное знание внешних воздействий.
\end{itemize}

Погрешность, связанную с неточностью исходной физико-математической
постановки, часто называют \textbf{неустранимой на этапе численного
решения}: её нельзя уменьшить простым сгущением сетки. Для этого нужна
более адекватная модель или более точные данные.

\subsection{Погрешность численного метода}

Точный оператор задачи заменяется приближённым. Если
\[
\|y_h-y\|\le Ch^p,
\]
то говорят о сходимости порядка $p$ относительно параметра $h$.
Для методов решения дифференциальных уравнений необходимо различать
\textbf{аппроксимацию}, \textbf{устойчивость} и \textbf{сходимость}.
Малый локальный дефект схемы сам по себе не гарантирует малую глобальную
ошибку: возмущения не должны неограниченно усиливаться алгоритмом.

Практически погрешность дискретизации оценивают, например, по серии расчётов
$h$, $h/2$, $h/4$ или с помощью экстраполяции Ричардсона.

\subsection{Абсолютная и относительная погрешности}

Если $a$ --- точное значение, а $\widehat a$ --- приближённое, то
\[
\Delta a=|\widehat a-a|
\]
называется абсолютной погрешностью, а при $a\ne0$
\[
\delta a=\frac{|\widehat a-a|}{|a|}
\]
--- относительной погрешностью.

Для функции $y=f(a_1,\ldots,a_n)$ при малых возмущениях входных данных
используется линеаризация
\[
\Delta y\approx
\sum_{i=1}^n
\frac{\partial f}{\partial a_i}\Delta a_i.
\]
Она показывает, какие параметры сильнее всего влияют на результат, и
лежит в основе анализа чувствительности.

\subsection{Обусловленность задачи и устойчивость алгоритма}

Важно различать свойства \textbf{задачи} и \textbf{алгоритма}.
Обусловленность характеризует чувствительность точного решения к малым
изменениям исходных данных. Для СЛАУ
\[
Ax=b
\]
в согласованной норме используется число обусловленности
\[
\boxed{\kappa(A)=\|A\|\,\|A^{-1}\|}.
\]
Большое $\kappa(A)$ означает, что даже точный алгоритм не может устранить
сильное влияние ошибок данных.

Численно устойчивый алгоритм, напротив, не должен существенно усиливать
малые погрешности, возникающие в ходе вычислений. Поэтому плохо
обусловленная задача и неустойчивый алгоритм --- разные причины потери
точности.

\subsection{Вычисления с плавающей точкой}

В ЭВМ вещественные числа хранятся с конечным числом значащих разрядов.
В нормализованной двоичной системе число имеет вид
\[
x=\pm m\,2^e,
\]
где мантисса $m$ и порядок $e$ принадлежат конечным множествам.
Следовательно, большинство вещественных чисел приходится округлять.

Для основных арифметических операций при отсутствии переполнения и
потери порядка стандартная модель округления имеет вид
\[
\boxed{
\operatorname{fl}(x\circ y)
=(x\circ y)(1+\delta),
\qquad |\delta|\le u,
}
\]
где $u$ --- единичная погрешность округления, а $\circ$ обозначает одну из
основных арифметических операций. Машинный эпсилон
$\varepsilon_{\rm mach}$ обычно определяют как расстояние от $1$ до
следующего представимого числа; при округлении к ближайшему $u$ имеет тот
же порядок и обычно равен половине этого расстояния.

\subsection{Накопление и потеря значащих цифр}

Малые ошибки отдельных операций могут накапливаться. Особенно опасно
\textbf{вычитание близких чисел}: ведущие значащие цифры сокращаются, и
относительная ошибка резко возрастает. Например, выражение
\[
\sqrt{x+1}-\sqrt{x}
\]
при больших $x$ лучше вычислять в эквивалентной форме
\[
\frac{1}{\sqrt{x+1}+\sqrt{x}}.
\]
Порядок выполнения операций также может влиять на погрешность, например при
суммировании большого числа величин сильно различающихся масштабов.

\subsection{Баланс погрешностей}

Уменьшение шага дискретизации не всегда улучшает окончательный результат.
В численном дифференцировании типична оценка
\[
E(h)\lesssim C_1h^p+\frac{C_2\varepsilon}{h^q},
\]
где первое слагаемое --- погрешность метода, а второе отражает усиление
ошибок исходных данных и округления. Поэтому существует оптимальный
диапазон $h$; брать шаг ``как можно меньше'' неправильно.

Практическое правило состоит в согласовании уровней ошибок: обычно нет
смысла решать дискретную задачу с точностью, намного превосходящей
неопределённость исходных данных и самой математической модели.

\subsection{Что важно сказать на экзамене}

Нужно уверенно различать три основных класса: \textbf{погрешность модели и
данных, погрешность метода, вычислительную погрешность}. Отдельно следует
различать обусловленность задачи и устойчивость алгоритма. Для вычислений
с плавающей точкой важно помнить модель относительного округления и
явление потери значащих цифр.

\newpage
